\subsection{Which plane cubics can carry a second design?}
\label{geo:cubic-classification}

A second design is highly isotropic, but its support need not fill
projective space. The following classification identifies exactly how
a cubic equation can coexist with that isotropy. Three structures enter
the proof essentially: a rank inequality for matrix pencils, the
moment-map description of multipartite entanglement, and the order-three
grading of the exceptional Lie algebra \(E_6\).

Let \(\omega=e^{2\pi i/3}\), and define the nine Hesse rays
\begin{equation}
 \mathcal H=\left\{
 \frac{(0,1,-\omega^j)}{\sqrt2},
 \frac{(1,0,-\omega^j)}{\sqrt2},
 \frac{(1,-\omega^j,0)}{\sqrt2}:j=0,1,2\right\},
 \qquad
 \mu_{\mathcal H}=\frac19\sum_{h\in\mathcal H}\delta_h.
 \label{geo:hesse-points}
\end{equation}
They form a symmetric informationally complete set:
\(\operatorname{tr}(P_hP_{h'})=1/4\) for \(h\ne h'\).
Their incidence geometry and linear dependencies are discussed in
\cite{geo_hesse2009,geo_dang2013}.

\begin{theorem}[Complete cubic-support classification]
\label{geo:cubic}
For a nonzero homogeneous ternary cubic \(F\), a projective second design
supported on \(\{F=0\}\) exists if and only if, after a unitary coordinate
change and multiplication by a nonzero scalar, \(F\) is either
\[
 z_0^3+z_1^3+z_2^3-\lambda z_0z_1z_2,\qquad \lambda^3\ne27,
\]
or \(z_0z_1z_2\).
In the smooth case the measure is uniquely \(\mu_{\mathcal H}\), in the
corresponding unitary coordinates. On the coordinate triangle every such
measure is
\[
 \mu=\frac13(\nu_0+\nu_1+\nu_2),
\]
where \(\nu_i\) is supported on \(z_i=0,\ |z_j|=|z_k|\), and its
first two phase Fourier coefficients vanish.
\end{theorem}

The conclusion is unitary, rather than merely projective: a general
projective transformation changes the prescribed Hermitian moments.
Nodal or cuspidal cubics, a line plus an irreducible conic, nonorthogonal
triangles, concurrent triples, and repeated components are excluded.

\begin{proof}
\emph{Step 1: a determinant with two balanced tensor factors.}
Every ternary cubic has a \(3\times3\) linear determinantal representation
\[
 F(z)=\det M(z),\qquad M(z)=\sum_{i=0}^2z_iM_i.
 \tag{a}
\]
For a smooth cubic, use the Hesse normal form
\cite[Lemma 1]{geo_hesse2009} and the Moore matrix
\[
 \mathsf M(a,z)=
 \begin{pmatrix}
 a_0z_0&a_1z_2&a_2z_1\\
 a_2z_2&a_0z_1&a_1z_0\\
 a_1z_1&a_2z_0&a_0z_2
 \end{pmatrix},
 \quad
 \det\mathsf M=S(a)T(z)-T(a)S(z),
\]
where \(S(z)=\sum_i z_i^3\), \(T(z)=z_0z_1z_2\). Choose \(a\) on the
cubic with \(T(a)\ne0\).
For a singular irreducible cubic, put a singular point at
\([0:0:1]\), so \(F=zQ_2(x,y)+Q_3(x,y)\). If \(Q_2\) has distinct
linear factors, normalize it to \(xy\) and absorb the terms divisible
by \(xy\) into \(z\), obtaining \(xyz+ax^3+dy^3\).
If \(Q_2\) is a square, normalize it to \(y^2\), absorb its multiples
into \(z\), then translate \(x\) to obtain \(x^3-y^2z\).
The two resulting determinants are
\[
 \begin{pmatrix}x&0&y\\0&y&x\\-dy&-ax&z\end{pmatrix},
 \qquad
 \begin{pmatrix}x&y&0\\0&x&z\\-y&0&x\end{pmatrix}.
\]
A reducible cubic is a line times a conic; represent the conic by a
\(2\times2\) linear determinant and add the line as a diagonal block.
This includes repeated factors.

Regard the coefficients of \(M\) as
\(\Psi\in(\mathbb C^3)^{\otimes3}\), and put
\[
 s=\|\Psi\|^2,\quad
 R=\sum_iM_iM_i^*,\quad
 C=\sum_iM_i^*M_i,\quad
 T_{ij}=\operatorname{tr}(M_iM_j^*).
\]
Up to a transpose for the second factor, these are its one-factor
reduced density matrices, each of trace \(s\).
The coefficients of \(\det M(z)\) are invariant under
\(M\mapsto AMB^{\mathsf T}\), \(A,B\in SL_3\).
Their nonvanishing prevents zero from belonging to the orbit closure.
The closure has a nonzero point of minimum Euclidean norm. Differentiating
its norm in all traceless Hermitian directions on the first two factors
gives
\[
 R=C=\frac{s}{3}I.
 \tag{b}
\]
The determinant remains exactly \(F\) at the limit. This argument does
not assume a finite invertible scaling of the original pencil.

\emph{Step 2: the design completes the balance.}
The Hermitian quartic
\[
 q_M(z)=2\operatorname{tr}\bigl((M(z)M(z)^*)^2\bigr)
       -\bigl(\operatorname{tr}(M(z)M(z)^*)\bigr)^2
\]
is nonnegative on \(\det M=0\). Indeed, a positive matrix of rank at most
two, with possibly zero eigenvalues \(b_1,b_2\), gives
\(q_M=(b_1-b_2)^2\). Haar's second-moment contractions give
\[
 12\int q_M\,d\sigma
 =2\operatorname{tr}R^2+2\operatorname{tr}C^2
      -\operatorname{tr}T^2-s^2.
 \tag{c}
\]
If a second design is supported on the cubic, then (b), (c), and its
moment identity imply
\[
 0\le12\int q_M\,d\mu
 =-\left\|T-\frac{s}{3}I\right\|_{\rm HS}^2\le0.
\]
Thus the third marginal is scalar. The measure constraint has forced a
three-qutrit tensor to be fully balanced.

\fig{fig03_tensor_balance}{The tensor used in the cubic classification.
The three matrix slices display the nonzero coefficients of a vector in
the Cartan subspace, with illustrative values $a=1$, $b=0.65$, and $c=0.35$.
After normalization, all three one-qutrit marginals are $I_3/3$, regardless
of these coefficients. The slice determinant is
$(a^3+b^3+c^3)z_0z_1z_2-abc(z_0^3+z_1^3+z_2^3)$, a member of the Hesse
pencil. The proof first obtains full marginal balance from the design
identity and only then invokes the $E_6$ Cartan representation.}{fig:tensor}

\emph{Step 3: the exceptional-Lie-algebra input.}
For the \(SL_3^3\) action, scalar one-factor marginals characterize
critical vectors. Their orbits are closed, and two critical vectors in
one orbit are locally unitarily equivalent, by Kempf--Ness
\cite[Theorem 14]{geo_sauerwein2018}. Vinberg's Cartan theorem for the
order-three \(E_6\) representation states that every closed orbit meets
\[
 \mathfrak c=\operatorname{span}\left\{
 |000\rangle+|111\rangle+|222\rangle,\,
 |012\rangle+|120\rangle+|201\rangle,\,
 |021\rangle+|210\rangle+|102\rangle\right\}.
\]
The statement includes nongeneric semisimple orbits
\cite[Section 7.1]{geo_ditrani2023}. Every vector in this space already
has scalar marginals: the three supports are disjoint Latin squares.
The tensor from (b) is consequently related to \(a u_1+b u_2+c u_3\)
by local unitaries. Its determinant is
\[
 (a^3+b^3+c^3)z_0z_1z_2
             -abc(z_0^3+z_1^3+z_2^3).
 \tag{d}
\]
The first two local unitaries change the determinant by a scalar; the
third gives a unitary change of its variables. Its nonzero singular
members are precisely orthogonal triangles. For example,
\[
 x^3+y^3+z^3-3xyz
 =(x+y+z)(x+\omega y+\omega^2z)(x+\omega^2y+\omega z),
\]
and the three coefficient vectors are orthogonal.

\emph{Step 4: rigidity on a smooth Hesse cubic.}
Define
\[
 r(z)=\sum_i|z_i|^4-\frac12\|z\|^4,\qquad
 D(z)=\sum_i z_i^2\overline{z_{i+1}z_{i+2}}.
\]
The Moore matrix satisfies
\[
 q_{\mathsf M(a,\cdot)}(z)
 =4r(a)r(z)+8\operatorname{Re}(D(a)\overline{D(z)}).
\]
On a common Hesse cubic, therefore, the vectors
\(v(z)=(r(z),\sqrt2\operatorname{Re}D(z),\sqrt2\operatorname{Im}D(z))\)
have pairwise nonnegative inner products. A second design has mean
\(v=0\). The nonnegative continuous kernel \(v(a)\cdot v(z)\) then has
zero double integral, forcing \(v=0\) on the entire support.

Suppose all coordinates of such a point are nonzero. For a unit
representative, put \(p_i=|z_i|^2\), \(u_i=z_i^3/p_i\).
The equations \(r=D=0\) imply
\[
 \sum_i u_i=0,\qquad |u_i|^2=p_i,\qquad \sum_i p_i^2=\frac12.
\]
Heron's identity gives zero area for the triangle of vectors \(u_i\).
Hence \(u_i=e^{i\theta}s_i\), where \(s_i\) are real,
\(\sum_i s_i=0\), and \(s_i^2=p_i\). Thus
\(z_i=e^{i\theta/3}\zeta_i s_i\), \(\zeta_i^3=1\), and
\[
 \frac{S(z)}{T(z)}=\frac3{\zeta_0\zeta_1\zeta_2}.
\]
This is a singular Hesse parameter. On a smooth member one coordinate
must therefore vanish; \(r=0\) equalizes the other two moduli, and the
cubic equation makes their cubes opposite. The support is contained
in \(\mathcal H\). Pairing \(M_2=I_6/6\) with \(P_i^{\otimes2}\)
gives
\[
 \frac16=w_i+\frac1{16}(1-w_i),
\]
so \(w_i=1/9\). Conversely those weights form a second design, by direct
summation or equality in \(\operatorname{tr}M_2^2\ge1/6\).

\emph{Step 5: the triangle.}
On \(xyz=0\), \(\sum_i|z_i|^4\ge1/2\), whose second-design expectation
is \(1/2\). Thus the support lies on the three midpoint circles.
The mixed moment \(\int|z_j|^2|z_k|^2\,d\mu=1/12\) gives mass \(1/3\)
to each circle. The off-diagonal identities are precisely vanishing
of its first and second phase Fourier coefficients. Conversely those
conditions give every bidegree-\((2,2)\) moment.
\end{proof}

\paragraph{Why this tensor format is special.}
For an \(n\times n\) pencil in \(d\) variables, the rank witness
\((n-1)\operatorname{tr}A^2-(\operatorname{tr}A)^2\) is nonnegative on
the determinant hypersurface. After balancing its first two factors,
a supported second design would imply
\[
 0\le\frac{n-2}{n}s^2-\operatorname{tr}T^2
 \le\left(\frac{n-2}{n}-\frac1d\right)s^2.
\]
The strictly negative range excludes designs. Equality forces the
third factor to balance, and for integers \(n,d\ge2\) occurs exactly at
\((n,d)=(3,3)\) and \((4,2)\). The binary-quartic case recovers a regular
tetrahedron on the Bloch sphere with uniform weights. Indeed, for a
second design on \(\mathbb P^1\), the nonnegative polynomial
\((3|\langle p,z\rangle|^2-1)^2\) has mean one and value four at \(p\),
so every atom has mass at most \(1/4\). A quartic has at most four
distinct zeros; equality therefore forces four equal weights, with
the polynomial zero at every other support point. Their squared
overlaps are \(1/3\), precisely the regular tetrahedron condition.
The three-qutrit case is a critical format of the rank inequality,
which explains the effectiveness of the exceptional-group description.
